\chapter{Tensor Products}
\label{ch:v10}

The tensor product converts bilinear maps into linear maps. It controls base change, localization, flatness, and many multilinear constructions throughout algebra and geometry.

\section*{Learning goals}
After this chapter, the reader should be able to:
\begin{itemize}
\item construct tensor products from bilinear maps using their universal property;
\item compute elementary tensor products of cyclic modules and quotients;
\item use tensor relations to simplify pure tensors and detect zero tensors;
\item prove right exactness of tensor product and identify where injectivity may fail;
\item apply base-change identities involving tensor products of algebras and modules;
\item distinguish tensor product from direct product, direct sum, and ordinary multiplication;
\end{itemize}
\section*{Conceptual roadmap}
\[
\boxed{\text{presentation or universal property}\;\longrightarrow\;\text{exactness/localization}\;\longrightarrow\;\text{structural consequence}.}
\]

\section{Balanced bilinear maps}
An $A$-balanced map satisfies $b(am,n)=b(m,an)$ in addition to additivity in each variable.

\section{Tensor product}
$M\otimes_A N$ is generated by symbols $m\otimes n$ modulo additivity and balancing relations.

% BEGIN VOL05-EXPANSION V10-example-01
\begin{example}[Cyclic tensor product]\label{ex:v10-expand-01}
For positive integers $m,n$, $\mathbb Z/m\otimes\mathbb Z/n\cong\mathbb Z/\gcd(m,n)$. Thus $\mathbb Z/6\otimes\mathbb Z/15\cong\mathbb Z/3$.
\end{example}
% END VOL05-EXPANSION V10-example-01
\section{Universal property}
Balanced bilinear maps out of $M\times N$ correspond uniquely to linear maps out of $M\otimes_A N$.

\section{Basic identities}
$A\otimes_A M\cong M$ and direct sums commute with tensor products.

\section{Quotient tensors}
$(A/I)\otimes_A M\cong M/IM$.

% BEGIN VOL05-EXPANSION V10-example-02
\begin{example}[Tensor with a quotient]\label{ex:v10-expand-02}
For an ideal $I$ and module $M$, $(R/I)\otimes_RM\cong M/IM$. Tensor relations force elements of $I$ to act as zero.
\end{example}
% END VOL05-EXPANSION V10-example-02
\section{Associativity and symmetry}
Canonical isomorphisms identify iterated tensor products and swap factors over a commutative ring.

\section{Right exactness}
Tensoring preserves cokernels and surjections, but not necessarily injections.

% BEGIN VOL05-EXPANSION V10-example-03
\begin{example}[Extension of scalars]\label{ex:v10-expand-03}
For a field extension $K/k$, $K\otimes_k k^n\cong K^n$. A $k$-basis becomes a $K$-basis after scalar extension.
\end{example}
% END VOL05-EXPANSION V10-example-03
\section{Localization}
$S^{-1}A\otimes_A M\cong S^{-1}M$.

\section{Examples}
Cyclic groups, quotient modules, field extensions, and polynomial base change exhibit torsion and dimension phenomena.

\section{Multilinear extension}
Exterior and symmetric powers are built by further quotienting tensor powers.

\section{Core structural results}

\begin{theorem}[Universal property]\label{thm:v10-01}
For every balanced bilinear map $b:M\times N\to P$, there is a unique linear map $\widetilde b:M\otimes_A N\to P$ with $\widetilde b(m\otimes n)=b(m,n)$.
\begin{proof}
Construct the tensor product as the free module on pairs modulo the subgroup generated by additivity and balancing relations. The relations are exactly those killed by every balanced bilinear map.
\end{proof}
\end{theorem}

\begin{theorem}[Quotient tensor identity]\label{thm:v10-02}
For an ideal $I\subset A$, $(A/I)\otimes_A M\cong M/IM$.
\begin{proof}
Send $(a+I)\otimes m$ to $am+IM$. The map is balanced and surjective; the universal properties on both sides show that its kernel is precisely the relations generated by $I M$.
\end{proof}
\end{theorem}

\begin{theorem}[Right exactness of tensor]\label{thm:v10-03}
If $M'\to M\to M''\to0$ is exact, then $M'\otimes N\to M\otimes N\to M''\otimes N\to0$ is exact.
\begin{proof}
The tensor product is a left adjoint in either variable, so it preserves cokernels; directly, generators in the kernel of the final map differ by tensors coming from the preceding module.
\end{proof}
\end{theorem}

\section{Worked examples}

\begin{example}[Balanced bilinear maps diagnostic]\label{ex:v10-worked-01}
Give a rigorous worked analysis of Balanced bilinear maps. State the ring, module, finiteness, localization, or homological hypotheses and derive the central conclusion. An $A$-balanced map satisfies $b(am,n)=b(m,an)$ in addition to additivity in each variable. The decisive step is to use the universal property, exactness criterion, localization test, or finite-generation mechanism appropriate to the algebraic structure.
\end{example}

\begin{example}[Tensor product diagnostic]\label{ex:v10-worked-02}
Give a rigorous worked analysis of Tensor product. State the ring, module, finiteness, localization, or homological hypotheses and derive the central conclusion. $M\otimes_A N$ is generated by symbols $m\otimes n$ modulo additivity and balancing relations. The decisive step is to use the universal property, exactness criterion, localization test, or finite-generation mechanism appropriate to the algebraic structure.
\end{example}

\begin{example}[Universal property diagnostic]\label{ex:v10-worked-03}
Give a rigorous worked analysis of Universal property. State the ring, module, finiteness, localization, or homological hypotheses and derive the central conclusion. Balanced bilinear maps out of $M\times N$ correspond uniquely to linear maps out of $M\otimes_A N$. The decisive step is to use the universal property, exactness criterion, localization test, or finite-generation mechanism appropriate to the algebraic structure.
\end{example}

\begin{example}[Basic identities diagnostic]\label{ex:v10-worked-04}
Give a rigorous worked analysis of Basic identities. State the ring, module, finiteness, localization, or homological hypotheses and derive the central conclusion. $A\otimes_A M\cong M$ and direct sums commute with tensor products. The decisive step is to use the universal property, exactness criterion, localization test, or finite-generation mechanism appropriate to the algebraic structure.
\end{example}

\section{Solved dossiers}
Each dossier is a canonical solved problem. The provenance ledger records which retained corpus rules guided its topic and which dossiers were added freshly to complete the chapter.

\begin{problem}[Balanced bilinear maps diagnostic]\label{prob:v10-dossier-01}
Give a rigorous worked analysis of Balanced bilinear maps. State the ring, module, finiteness, localization, or homological hypotheses and derive the central conclusion.
\end{problem}
\begin{solution}
An $A$-balanced map satisfies $b(am,n)=b(m,an)$ in addition to additivity in each variable. The decisive step is to use the universal property, exactness criterion, localization test, or finite-generation mechanism appropriate to the algebraic structure.
\end{solution}

\begin{problem}[Tensor product diagnostic]\label{prob:v10-dossier-02}
Give a rigorous worked analysis of Tensor product. State the ring, module, finiteness, localization, or homological hypotheses and derive the central conclusion.
\end{problem}
\begin{solution}
$M\otimes_A N$ is generated by symbols $m\otimes n$ modulo additivity and balancing relations. The decisive step is to use the universal property, exactness criterion, localization test, or finite-generation mechanism appropriate to the algebraic structure.
\end{solution}

\begin{problem}[Universal property diagnostic]\label{prob:v10-dossier-03}
Give a rigorous worked analysis of Universal property. State the ring, module, finiteness, localization, or homological hypotheses and derive the central conclusion.
\end{problem}
\begin{solution}
Balanced bilinear maps out of $M\times N$ correspond uniquely to linear maps out of $M\otimes_A N$. The decisive step is to use the universal property, exactness criterion, localization test, or finite-generation mechanism appropriate to the algebraic structure.
\end{solution}

\begin{problem}[Basic identities diagnostic]\label{prob:v10-dossier-04}
Give a rigorous worked analysis of Basic identities. State the ring, module, finiteness, localization, or homological hypotheses and derive the central conclusion.
\end{problem}
\begin{solution}
$A\otimes_A M\cong M$ and direct sums commute with tensor products. The decisive step is to use the universal property, exactness criterion, localization test, or finite-generation mechanism appropriate to the algebraic structure.
\end{solution}

\begin{problem}[Quotient tensors diagnostic]\label{prob:v10-dossier-05}
Give a rigorous worked analysis of Quotient tensors. State the ring, module, finiteness, localization, or homological hypotheses and derive the central conclusion.
\end{problem}
\begin{solution}
$(A/I)\otimes_A M\cong M/IM$. The decisive step is to use the universal property, exactness criterion, localization test, or finite-generation mechanism appropriate to the algebraic structure.
\end{solution}

\begin{problem}[Associativity and symmetry diagnostic]\label{prob:v10-dossier-06}
Give a rigorous worked analysis of Associativity and symmetry. State the ring, module, finiteness, localization, or homological hypotheses and derive the central conclusion.
\end{problem}
\begin{solution}
Canonical isomorphisms identify iterated tensor products and swap factors over a commutative ring. The decisive step is to use the universal property, exactness criterion, localization test, or finite-generation mechanism appropriate to the algebraic structure.
\end{solution}

\begin{problem}[Right exactness diagnostic]\label{prob:v10-dossier-07}
Give a rigorous worked analysis of Right exactness. State the ring, module, finiteness, localization, or homological hypotheses and derive the central conclusion.
\end{problem}
\begin{solution}
Tensoring preserves cokernels and surjections, but not necessarily injections. The decisive step is to use the universal property, exactness criterion, localization test, or finite-generation mechanism appropriate to the algebraic structure.
\end{solution}

\begin{problem}[Localization diagnostic]\label{prob:v10-dossier-08}
Give a rigorous worked analysis of Localization. State the ring, module, finiteness, localization, or homological hypotheses and derive the central conclusion.
\end{problem}
\begin{solution}
$S^{-1}A\otimes_A M\cong S^{-1}M$. The decisive step is to use the universal property, exactness criterion, localization test, or finite-generation mechanism appropriate to the algebraic structure.
\end{solution}

\begin{problem}[Universal property proof dossier]\label{prob:v10-dossier-09}
Prove the following structural statement and identify the hypothesis that makes the conclusion work: For every balanced bilinear map $b:M\times N\to P$, there is a unique linear map $\widetilde b:M\otimes_A N\to P$ with $\widetilde b(m\otimes n)=b(m,n)$.
\end{problem}
\begin{solution}
Construct the tensor product as the free module on pairs modulo the subgroup generated by additivity and balancing relations. The relations are exactly those killed by every balanced bilinear map. This proof isolates the reusable algebraic mechanism behind the result.
\end{solution}

\begin{problem}[Quotient tensor identity proof dossier]\label{prob:v10-dossier-10}
Prove the following structural statement and identify the hypothesis that makes the conclusion work: For an ideal $I\subset A$, $(A/I)\otimes_A M\cong M/IM$.
\end{problem}
\begin{solution}
Send $(a+I)\otimes m$ to $am+IM$. The map is balanced and surjective; the universal properties on both sides show that its kernel is precisely the relations generated by $I M$. This proof isolates the reusable algebraic mechanism behind the result.
\end{solution}

\begin{problem}[Right exactness of tensor proof dossier]\label{prob:v10-dossier-11}
Prove the following structural statement and identify the hypothesis that makes the conclusion work: If $M'\to M\to M''\to0$ is exact, then $M'\otimes N\to M\otimes N\to M''\otimes N\to0$ is exact.
\end{problem}
\begin{solution}
The tensor product is a left adjoint in either variable, so it preserves cokernels; directly, generators in the kernel of the final map differ by tensors coming from the preceding module. This proof isolates the reusable algebraic mechanism behind the result.
\end{solution}

\begin{problem}[Chapter synthesis]\label{prob:v10-dossier-12}
Connect the definitions and principal structural theorems of this chapter into one reusable commutative-algebra strategy.
\end{problem}
\begin{solution}
The tensor product converts bilinear maps into linear maps. It controls base change, localization, flatness, and many multilinear constructions throughout algebra and geometry. The recurring strategy is to encode the algebraic object by kernels, quotients, localization, finite presentation, or a resolution, apply the corresponding universal or exactness principle, and then recover the desired global or local conclusion.
\end{solution}

\section{Exercises with complete solutions}

\begin{exercise}\label{exr:v10-01}
State and apply the central principle behind Balanced bilinear maps. Give one concrete algebraic consequence.
\end{exercise}
\begin{hint}
Use the universal property: a bilinear map M x N -> P factors uniquely through M tensor N.
\end{hint}
\begin{solution}
An $A$-balanced map satisfies $b(am,n)=b(m,an)$ in addition to additivity in each variable. The same principle can then be reused in later localization, support, base-change, resolution, Tor, or Ext arguments.
\end{solution}

\begin{exercise}\label{exr:v10-02}
State and apply the central principle behind Tensor product. Give one concrete algebraic consequence.
\end{exercise}
\begin{hint}
Apply relations (m+m') tensor n, m tensor (n+n'), and (am) tensor n = m tensor (an) to simplify tensors.
\end{hint}
\begin{solution}
$M\otimes_A N$ is generated by symbols $m\otimes n$ modulo additivity and balancing relations. The same principle can then be reused in later localization, support, base-change, resolution, Tor, or Ext arguments.
\end{solution}

\begin{exercise}\label{exr:v10-03}
State and apply the central principle behind Universal property. Give one concrete algebraic consequence.
\end{exercise}
\begin{hint}
For cyclic modules, start from \(A/I\otimes_A N\cong N/IN\) and reduce the computation to a quotient.
\end{hint}
\begin{solution}
Balanced bilinear maps out of $M\times N$ correspond uniquely to linear maps out of $M\otimes_A N$. The same principle can then be reused in later localization, support, base-change, resolution, Tor, or Ext arguments.
\end{solution}

\begin{exercise}\label{exr:v10-04}
State and apply the central principle behind Basic identities. Give one concrete algebraic consequence.
\end{exercise}
\begin{hint}
To prove right exactness, tensor a presentation and identify the resulting cokernel.
\end{hint}
\begin{solution}
$A\otimes_A M\cong M$ and direct sums commute with tensor products. The same principle can then be reused in later localization, support, base-change, resolution, Tor, or Ext arguments.
\end{solution}

\begin{exercise}\label{exr:v10-05}
State and apply the central principle behind Quotient tensors. Give one concrete algebraic consequence.
\end{exercise}
\begin{hint}
Injectivity can fail after tensoring; look for torsion killed by multiplication before and after tensoring.
\end{hint}
\begin{solution}
$(A/I)\otimes_A M\cong M/IM$. The same principle can then be reused in later localization, support, base-change, resolution, Tor, or Ext arguments.
\end{solution}

\begin{exercise}\label{exr:v10-06}
State and apply the central principle behind Associativity and symmetry. Give one concrete algebraic consequence.
\end{exercise}
\begin{hint}
Use pure tensors only as generators; not every tensor is itself pure.
\end{hint}
\begin{solution}
Canonical isomorphisms identify iterated tensor products and swap factors over a commutative ring. The same principle can then be reused in later localization, support, base-change, resolution, Tor, or Ext arguments.
\end{solution}

\begin{exercise}\label{exr:v10-07}
State and apply the central principle behind Right exactness. Give one concrete algebraic consequence.
\end{exercise}
\begin{hint}
For base change, reassociate tensors carefully and keep track of which ring acts at each stage.
\end{hint}
\begin{solution}
Tensoring preserves cokernels and surjections, but not necessarily injections. The same principle can then be reused in later localization, support, base-change, resolution, Tor, or Ext arguments.
\end{solution}

\begin{exercise}\label{exr:v10-08}
State and apply the central principle behind Localization. Give one concrete algebraic consequence.
\end{exercise}
\begin{hint}
Distinguish direct sum from tensor product by their universal properties: one is additive coproduct, the other linearizes bilinear maps.
\end{hint}
\begin{solution}
$S^{-1}A\otimes_A M\cong S^{-1}M$. The same principle can then be reused in later localization, support, base-change, resolution, Tor, or Ext arguments.
\end{solution}

\section*{Chapter summary}
The chapter now contributes a canonical solved layer to the progression from rings and modules through localization, Noetherian methods, integral dependence, and derived functors.
% BEGIN VOL05-EXPANSION V10-exercises-01
\section{Graded supplementary exercises}
These exercises extend the protected chapter with a balanced set of computations, proofs, hypothesis tests, applications, and challenges.

\subsection*{Standard computations}

\begin{exercise}[Unit tensor]\label{exr:v10-expand-01}
Compute $R\otimes_RM$.
\end{exercise}
\begin{hint}
Use tensor products and their universal property.
\end{hint}
\begin{solution}
It is $M$.
\end{solution}

\begin{exercise}[Zero tensor]\label{exr:v10-expand-02}
Compute $M\otimes_R0$.
\end{exercise}
\begin{hint}
Use tensor products and their universal property.
\end{hint}
\begin{solution}
It is zero.
\end{solution}

\begin{exercise}[Cyclic tensor]\label{exr:v10-expand-03}
Compute $\mathbb Z/4\otimes\mathbb Z/6$.
\end{exercise}
\begin{hint}
Use tensor products and their universal property.
\end{hint}
\begin{solution}
It is $\mathbb Z/2$.
\end{solution}

\begin{exercise}[Coprime tensor]\label{exr:v10-expand-04}
Compute $\mathbb Z/4\otimes\mathbb Z/9$.
\end{exercise}
\begin{hint}
Use tensor products and their universal property.
\end{hint}
\begin{solution}
It is zero.
\end{solution}

\begin{exercise}[Quotient tensor]\label{exr:v10-expand-05}
Compute $(R/I)\otimes_R(R/J)$.
\end{exercise}
\begin{hint}
Use tensor products and their universal property.
\end{hint}
\begin{solution}
It is $R/(I+J)$.
\end{solution}

\subsection*{Proofs}

\begin{exercise}[Universal property proof]\label{exr:v10-expand-06}
Prove: For every balanced bilinear map $b:M\times N\to P$, there is a unique linear map $\widetilde b:M\otimes_A N\to P$ with $\widetilde b(m\otimes n)=b(m,n)$.
\end{exercise}
\begin{hint}
Start from the chapter theorem hypotheses and identify the decisive algebraic map or containment.
\end{hint}
\begin{solution}
Construct the tensor product as the free module on pairs modulo the subgroup generated by additivity and balancing relations. The relations are exactly those killed by every balanced bilinear map.
\end{solution}

\begin{exercise}[Quotient tensor identity proof]\label{exr:v10-expand-07}
Prove: For an ideal $I\subset A$, $(A/I)\otimes_A M\cong M/IM$.
\end{exercise}
\begin{hint}
Start from the chapter theorem hypotheses and identify the decisive algebraic map or containment.
\end{hint}
\begin{solution}
Send $(a+I)\otimes m$ to $am+IM$. The map is balanced and surjective; the universal properties on both sides show that its kernel is precisely the relations generated by $I M$.
\end{solution}

\begin{exercise}[Right exactness of tensor proof]\label{exr:v10-expand-08}
Prove: If $M'\to M\to M''\to0$ is exact, then $M'\otimes N\to M\otimes N\to M''\otimes N\to0$ is exact.
\end{exercise}
\begin{hint}
Start from the chapter theorem hypotheses and identify the decisive algebraic map or containment.
\end{hint}
\begin{solution}
The tensor product is a left adjoint in either variable, so it preserves cokernels; directly, generators in the kernel of the final map differ by tensors coming from the preceding module.
\end{solution}

\begin{exercise}[Structural synthesis]\label{exr:v10-expand-09}
Prove a short corollary by combining Universal property with Quotient tensor identity. State every hypothesis you use.
\end{exercise}
\begin{hint}
Apply the first theorem to move to its structural description, then use the second theorem on that description.
\end{hint}
\begin{solution}
A valid synthesis first invokes Universal property and then applies Quotient tensor identity; the conclusion follows after checking the shared hypotheses.
\end{solution}

\subsection*{Counterexamples and hypothesis tests}

\begin{exercise}[Left exactness]\label{exr:v10-expand-10}
Test the claim: Tensoring always preserves injections.
\end{exercise}
\begin{hint}
Try a smallest standard ring or module from this chapter.
\end{hint}
\begin{solution}
False: tensor multiplication by $2$ with $\mathbb Z/2$.
\end{solution}

\begin{exercise}[Pure tensor]\label{exr:v10-expand-11}
Test the claim: A pure tensor is zero only if one factor is zero.
\end{exercise}
\begin{hint}
Try a smallest standard ring or module from this chapter.
\end{hint}
\begin{solution}
False: use $\mathbb Z/2\otimes\mathbb Z/3=0$.
\end{solution}

\begin{exercise}[Dimension rule]\label{exr:v10-expand-12}
Test the claim: Dimension multiplication holds for arbitrary modules.
\end{exercise}
\begin{hint}
Try a smallest standard ring or module from this chapter.
\end{hint}
\begin{solution}
False; it is a vector-space formula under suitable finiteness.
\end{solution}

\subsection*{Applications and investigations}

\begin{exercise}[Scalar extension]\label{exr:v10-expand-13}
Why is $B\otimes_AM$ extension of scalars?
\end{exercise}
\begin{hint}
Translate the situation into tensor products and their universal property.
\end{hint}
\begin{solution}
It is universal for reinterpreting the $A$-action through $B$.
\end{solution}

\begin{exercise}[Affine products]\label{exr:v10-expand-14}
Why do tensor products of algebras model product coordinates?
\end{exercise}
\begin{hint}
Translate the situation into tensor products and their universal property.
\end{hint}
\begin{solution}
They combine independent coordinate functions over the base field.
\end{solution}

\subsection*{Challenge problems}

\begin{exercise}[Local-global synthesis]\label{exr:v10-expand-15}
Connect the chapter ideas 'Balanced bilinear maps' and 'Multilinear extension' in one rigorous argument.
\end{exercise}
\begin{hint}
State the relevant map, ideal, module, or universal property before drawing the conclusion.
\end{hint}
\begin{solution}
The argument begins with An $A$-balanced map satisfies $b(am,n)=b(m,an)$ in addition to additivity in each variable. It then uses the later viewpoint: Exterior and symmetric powers are built by further quotienting tensor powers. The bridge is supplied by the structural theorems proved in the chapter.
\end{solution}

\begin{exercise}[Hypothesis audit]\label{exr:v10-expand-16}
Choose one theorem from the chapter and explain precisely what can fail if one key hypothesis is removed.
\end{exercise}
\begin{hint}
Use one of the explicit counterexamples in the graded set and compare it with the theorem statement.
\end{hint}
\begin{solution}
The theorem hypotheses are essential because the chapter's quotient, localization, exactness, or finiteness mechanism can fail outside them. A correct answer identifies the dropped hypothesis and exhibits a concrete failure.
\end{solution}

% END VOL05-EXPANSION V10-exercises-01
